% !TEX root = ../main.tex
% File: chapters_part1/chap4_5.tex
% Nội dung cho Chương 4, Phần 5

\section{Lý thuyết về Mixture-of-Experts (MoE)}
\label{sec:mixture_of_experts}

Khi chúng ta nói về việc “tăng quy mô” (scaling up) các mô hình ngôn ngữ, cách tiếp cận thông thường là làm cho các mô hình trở nên “dày đặc” hơn (denser) - tức là tăng số lớp, tăng số chiều ẩn, tăng số đầu chú ý. Tuy nhiên, cách tiếp cận này có một giới hạn: để xử lý mỗi token, \textbf{toàn bộ các tham số của mô hình đều phải được kích hoạt}. Điều này dẫn đến chi phí tính toán cực kỳ lớn ở cả quá trình huấn luyện và suy luận.

Mixture-of-Experts (MoE) \cite{shazeer2017outrageously, fedus2022switch} đưa ra một giải pháp thanh lịch cho vấn đề này, dựa trên một triết lý rất con người:

\begin{tcolorbox}[
    title=Triết lý của Mixture-of-Experts,
    colback=yellow!10!white, colframe=yellow!50!black, fonttitle=\bfseries
]
“Không cần phải huy động toàn bộ bộ não để giải quyết mọi vấn đề. Thay vào đó, hãy có một hội đồng gồm nhiều \textbf{chuyên gia (experts)}, mỗi người chuyên về một lĩnh vực khác nhau. Đối với mỗi vấn đề cụ thể, chỉ cần hỏi ý kiến của một vài chuyên gia phù hợp nhất.”
\end{tcolorbox}

Trong bối cảnh của Transformer, MoE cho phép tăng số lượng tham số của mô hình lên rất lớn, nhưng lại giữ cho chi phí tính toán trên mỗi token không đổi hoặc chỉ tăng nhẹ. Đây được gọi là \textbf{tăng quy mô thưa thớt (sparse scaling)}.

\subsection{Kiến trúc và Nguyên lý hoạt động}
\label{ssec:moe_architecture}

Kiến trúc MoE không phải là một mô hình hoàn toàn mới. Thay vào đó, nó là một cách để sửa đổi các khối Transformer hiện có. Cụ thể, trong một khối Transformer tiêu chuẩn, lớp \textbf{Feed-Forward Network (FFN)} được thay thế bằng một \textbf{lớp MoE}.

Một lớp MoE bao gồm hai thành phần chính:
\begin{enumerate}
    \item \textbf{Một tập hợp gồm $N$ mạng FFN độc lập, gọi là các “chuyên gia” (Experts).} Ví dụ, $N$ có thể là 8, 16, hoặc 64. Mỗi chuyên gia này có bộ trọng số riêng.
    \item \textbf{Một mạng cổng (Gating Network) nhỏ.} Đây là một mạng nơ-ron có nhiệm vụ quyết định xem nên gửi mỗi token đến chuyên gia nào.
\end{enumerate}

\begin{center}
    \includegraphics[width=1.0\textwidth]{moe_transformer_block.png}
    \captionof{figure}{So sánh một khối Transformer tiêu chuẩn (trái) và một khối Transformer với lớp MoE (phải). Lớp FFN dày đặc được thay thế bằng một mạng cổng và nhiều chuyên gia FFN thưa thớt.}
    \label{fig:moe_transformer_block}
\end{center}

\subsubsection{Dòng chảy dữ liệu của một token}
Khi một token (được biểu diễn bằng vector $x$) đi vào lớp MoE, quá trình sau sẽ xảy ra:

\paragraph{Bước 1: Định tuyến bởi Mạng cổng (Routing by the Gating Network)}
\begin{itemize}
    \item Token $x$ được đưa vào mạng cổng.
    \item Mạng cổng, thường là một lớp tuyến tính đơn giản, sẽ xuất ra một vector logits có số chiều bằng số chuyên gia ($N$).
    \item Một hàm `softmax` được áp dụng lên vector logits này để tạo ra một phân phối xác suất, $G(x) = (g_1, g_2, \dots, g_N)$, trong đó $g_i$ là “điểm số” hoặc “mức độ tin cậy” mà mạng cổng dành cho chuyên gia thứ $i$.
\end{itemize}

\paragraph{Bước 2: Chọn các Chuyên gia “Top-K”}
\begin{itemize}
    \item Thay vì gửi token đến tất cả các chuyên gia (điều này sẽ làm mất đi tính hiệu quả), hệ thống chỉ chọn ra $k$ chuyên gia có điểm số cao nhất (Top-K routing). Thông thường, $k$ là một số rất nhỏ, ví dụ $k=1$ hoặc $k=2$.
    \item Ví dụ, nếu $k=2$ và chuyên gia thứ 3 và thứ 5 có điểm số cao nhất, thì token này sẽ chỉ được gửi đến hai chuyên gia này. Tất cả các chuyên gia khác sẽ không được kích hoạt cho token này.
\end{itemize}

\paragraph{Bước 3: Xử lý bởi các Chuyên gia}
\begin{itemize}
    \item Token $x$ được đưa qua các mạng FFN của $k$ chuyên gia đã được chọn.
    \item Kết quả đầu ra là $k$ vector: $E_1(x), E_2(x), \dots, E_k(x)$.
\end{itemize}

\paragraph{Bước 4: Tổng hợp có trọng số}
\begin{itemize}
    \item Đầu ra cuối cùng của lớp MoE, $y$, là một tổng có trọng số của đầu ra từ các chuyên gia được chọn. Các trọng số này chính là các điểm số $g_i$ từ mạng cổng đã được chuẩn hóa.
    $$ y = \sum_{i \in \text{TopK}} g_i \cdot E_i(x) $$
\end{itemize}

\subsubsection{Ví dụ: Mixtral 8x7B}
Mô hình Mixtral 8x7B của Mistral AI \cite{jiang2024mixtral} là một ví dụ điển hình của MoE.
\begin{itemize}
    \item \textbf{8x7B có nghĩa là gì?} Mô hình được xây dựng từ kiến trúc Mistral 7B, trong đó FFN ở \emph{mỗi lớp} được thay bằng 8 chuyên gia. Một lầm tưởng phổ biến là tổng số tham số bằng $8 \times 7B = 56B$. Thực tế, chỉ các FFN được nhân bản; các lớp attention, embedding và LayerNorm vẫn \textbf{dùng chung}. Vì vậy tổng số tham số là khoảng \textbf{46.7B}.
    \item \textbf{Top-2 Routing:} Tại mỗi lớp MoE, mỗi token chỉ được định tuyến đến 2 trong số 8 chuyên gia.
    \item \textbf{Chi phí tính toán:} Mỗi token chỉ kích hoạt khoảng \textbf{12.9B} tham số (phần dùng chung cộng hai chuyên gia), nên chi phí suy luận trên mỗi token tương đương một mô hình dày đặc cỡ 13B, trong khi chất lượng ngang hoặc vượt Llama 2 70B trên nhiều benchmark. Lưu ý rằng \emph{bộ nhớ} vẫn phải chứa đủ 46.7B tham số.
\end{itemize}
\begin{tcolorbox}[
    title={Các “Chuyên gia” có thực sự chuyên môn hóa không?},
    colback=blue!5!white, colframe=blue!75!black, fonttitle=\bfseries
]
Phép ẩn dụ về “hội đồng chuyên gia” rất hữu ích, nhưng sự chuyên môn hóa của các expert FFN trong thực tế lại trừu tượng hơn nhiều.
\begin{itemize}
    \item \textbf{Không phải là chuyên môn về chủ đề:} Các nghiên cứu đã chỉ ra rằng các expert thường không chuyên môn hóa theo các chủ đề dễ hiểu như “thể thao”, “khoa học” hay “lịch sử”. Việc định tuyến một token đến expert nào phụ thuộc nhiều vào các đặc trưng trừu tượng của vector biểu diễn của nó.
    \item \textbf{Chuyên môn hóa theo cụm cú pháp/ngữ nghĩa:} Thay vào đó, sự chuyên môn hóa có thể xảy ra ở mức độ thấp hơn. Ví dụ, một số chuyên gia có thể giỏi hơn trong việc xử lý các token là dấu câu, một số khác chuyên về các động từ, một số khác lại tập trung vào các token ở các lớp Transformer cụ thể (ví dụ, các token ở lớp đầu so với lớp cuối).
    \item \textbf{Sự dư thừa có chủ đích:} Có bằng chứng cho thấy các chuyên gia không hoàn toàn độc lập. Thường có sự dư thừa (redundancy) giữa chúng. Điều này có thể là một đặc tính hữu ích, giúp mô hình trở nên mạnh mẽ hơn: nếu một chuyên gia đưa ra kết quả kém, các chuyên gia khác có thể bù đắp.
\end{itemize}
Vì vậy, hãy xem các “chuyên gia” như các \textbf{nhóm tham số (parameter groups)} có thể học được, thay vì các chuyên gia về một lĩnh vực kiến thức cụ thể. Mạng cổng học cách kết hợp các nhóm tham số này một cách tối ưu cho từng token.
\end{tcolorbox}
\subsection{Thách thức trong việc Huấn luyện MoE}
\label{ssec:moe_training_challenges}

Mặc dù rất mạnh mẽ, việc huấn luyện MoE không hề đơn giản và đòi hỏi phải giải quyết một số thách thức đặc thù.

\paragraph{1. Thách thức về sự sụp đổ của việc định tuyến (Routing Collapse)}
Đây là vấn đề lớn nhất của MoE. Nếu không được kiểm soát, mạng cổng sẽ nhanh chóng học được một chiến lược “lười biếng”: luôn gửi phần lớn các token đến một hoặc một vài chuyên gia “yêu thích” mà nó cho là tốt nhất.
\begin{itemize}
    \item \textbf{Hậu quả:}
        \begin{enumerate}
            \item \textbf{Lãng phí tham số:} Các chuyên gia khác không nhận được token, do đó không được huấn luyện và trở thành “dead wood”. Mô hình MoE khổng lồ sẽ sụp đổ (collapse) thành một mô hình nhỏ hơn nhiều.
            \item \textbf{Mất ổn định huấn luyện:} Gradient sẽ chảy không đều, gây khó khăn cho quá trình tối ưu hóa.
        \end{enumerate}
\end{itemize}

\paragraph{Giải pháp: Hàm Mất mát Cân bằng Tải (Load Balancing Loss)}
Để giải quyết vấn đề này, người ta thêm một \textbf{hàm mất mát phụ trợ (auxiliary loss)} vào hàm mất mát chính của mô hình. Hàm mất mát này được thiết kế để khuyến khích sự phân bổ token đồng đều và thường bao gồm hai thành phần:
\begin{itemize}
    \item \textbf{Một thành phần khuyến khích tầm quan trọng (importance) đồng đều:} Nó phạt nếu tổng các điểm số (logits) mà mạng cổng gán cho mỗi chuyên gia trên toàn bộ một batch là không đồng đều. Điều này đảm bảo mỗi chuyên gia được coi là “quan trọng” như nhau về mặt tổng thể.
    \item \textbf{Một thành phần khuyến khích số lượng token (load) đồng đều:} Nó phạt nếu số lượng token được gửi đến mỗi chuyên gia trong một batch là không đồng đều. Điều này đảm bảo mỗi chuyên gia có “công việc” để làm.
\end{itemize}
Việc cân bằng giữa hàm mất mát chính (ví dụ, cross-entropy) và hàm mất mát phụ trợ này là một trong những khía cạnh nghệ thuật nhất khi huấn luyện các mô hình MoE.

\subsection{Switch Transformer: Định tuyến Top-1 và Quy mô Nghìn tỷ Tham số \newtag}
\label{ssec:switch_transformer}
\textbf{Switch Transformer} \cite{fedus2022switch} đặt câu hỏi ngược với trực giác thời đó: \emph{có cần gửi mỗi token tới ít nhất hai chuyên gia không?} Câu trả lời là không -- định tuyến tới \textbf{đúng một} chuyên gia (top-1, gọi là “switch”) vẫn giữ được chất lượng, đồng thời giảm một nửa chi phí tính toán và giao tiếp của lớp MoE.

\paragraph{Dung lượng chuyên gia (expert capacity).} Để tính toán song song với tensor kích thước cố định, mỗi chuyên gia chỉ nhận tối đa
\[
\text{capacity} = \left(\frac{\text{số token trong batch}}{\text{số chuyên gia}}\right) \times \text{capacity factor}
\]
token. Token vượt dung lượng bị “tràn” (dropped) và chỉ đi qua kết nối phần dư. Capacity factor lớn hơn 1 (ví dụ 1.25) giảm token bị tràn nhưng tốn thêm bộ nhớ và giao tiếp.

\paragraph{Hàm mất mát cân bằng tải khả vi.} Với $N$ chuyên gia và một batch $\mathcal{B}$ gồm $T$ token:
\[
\mathcal{L}_{\text{aux}} = \alpha \cdot N \cdot \sum_{i=1}^{N} f_i \cdot P_i,
\qquad
f_i = \frac{1}{T}\sum_{x\in\mathcal{B}} \mathbb{1}\{\arg\max p(x) = i\},
\qquad
P_i = \frac{1}{T}\sum_{x\in\mathcal{B}} p_i(x),
\]
trong đó $f_i$ là tỷ lệ token thực sự được gửi tới chuyên gia $i$ (không khả vi) và $P_i$ là xác suất định tuyến trung bình dành cho chuyên gia $i$ (khả vi). Tổng $\sum f_i P_i$ đạt cực tiểu khi phân bổ đều ($f_i = P_i = 1/N$); hệ số $N$ giữ cho giá trị loss không đổi khi thay đổi số chuyên gia, và $\alpha = 10^{-2}$ trong bài báo.

\paragraph{Các “mẹo” ổn định huấn luyện.} (1) Tính router ở độ chính xác float32 trong khi phần còn lại dùng bfloat16 (\emph{selective precision}); (2) khởi tạo trọng số nhỏ hơn (giảm hệ số tỷ lệ khởi tạo xuống 10 lần); (3) dropout lớn hơn ở các lớp chuyên gia khi tinh chỉnh. Kết quả: Switch-Base đạt chất lượng của T5-Base nhanh hơn tới 7 lần với cùng FLOPs mỗi token, và mô hình Switch-C đạt 1.6 nghìn tỷ tham số với 2048 chuyên gia.

\subsection{DeepSeekMoE: Chuyên gia Mịn, Chuyên gia Chia sẻ và Cân bằng Tải không cần Loss phụ \newtag}
\label{ssec:deepseekmoe}
Các MoE kiểu GShard/Switch/Mixtral có ít chuyên gia (8--64) và mỗi chuyên gia khá lớn. DeepSeek chỉ ra hai vấn đề: \textbf{kiến thức bị trộn lẫn} (mỗi chuyên gia phải học quá nhiều thứ khác nhau) và \textbf{kiến thức bị dư thừa} (nhiều chuyên gia cùng phải học lại những kiến thức chung). \textbf{DeepSeekMoE} \cite{dai2024deepseekmoe} đề xuất hai thay đổi:
\begin{enumerate}
	\item \textbf{Phân đoạn chuyên gia mịn (fine-grained expert segmentation):} Chia mỗi chuyên gia thành $m$ chuyên gia nhỏ hơn (kích thước ẩn FFN giảm $m$ lần) và kích hoạt $mK$ chuyên gia thay vì $K$. Chi phí tính toán giữ nguyên, nhưng số tổ hợp chuyên gia có thể có tăng bùng nổ, cho phép chuyên môn hóa tinh tế hơn.
	\item \textbf{Cô lập chuyên gia chia sẻ (shared expert isolation):} Dành $K_s$ chuyên gia \emph{luôn được kích hoạt} cho mọi token để học kiến thức chung, giải phóng các chuyên gia định tuyến khỏi việc học trùng lặp.
\end{enumerate}
Đầu ra của lớp MoE cho token $u_t$ khi đó là
\[
h'_t = u_t + \sum_{i=1}^{N_s} \text{FFN}^{(s)}_i(u_t) + \sum_{i=1}^{N_r} g_{i,t}\,\text{FFN}^{(r)}_i(u_t),
\]
với $g_{i,t}$ khác 0 chỉ khi chuyên gia định tuyến $i$ nằm trong top-$K_r$ của token $t$.

\paragraph{Quy mô thực tế.} DeepSeek-V2 \cite{deepseekai2024v2} có 236B tham số (21B kích hoạt/token). DeepSeek-V3 \cite{deepseekai2024v3} có 671B tham số (37B kích hoạt/token); mỗi lớp MoE gồm 1 chuyên gia chia sẻ và 256 chuyên gia định tuyến, mỗi token chọn 8 chuyên gia, và bị giới hạn chỉ gửi tới tối đa 4 máy (node-limited routing) để kiểm soát chi phí giao tiếp.

\paragraph{Cân bằng tải không dùng loss phụ (auxiliary-loss-free).} Loss phụ trợ có một cái giá: nó tạo ra gradient “gây nhiễu” cạnh tranh với mục tiêu mô hình ngôn ngữ; hệ số $\alpha$ lớn thì cân bằng tốt nhưng chất lượng giảm. DeepSeek-V3 \cite{wang2024auxlossfree, deepseekai2024v3} bỏ hẳn loss phụ và thay bằng một \textbf{bias} $b_i$ cho mỗi chuyên gia, chỉ dùng khi \emph{chọn} top-$K$:
\[
g'_{i,t} = \begin{cases} s_{i,t} & \text{nếu } s_{i,t} + b_i \in \text{Top}_K\big(\{s_{j,t} + b_j\}_{j=1}^{N_r}\big) \\ 0 & \text{ngược lại}\end{cases}
\]
Trọng số trộn vẫn lấy từ điểm gốc $s_{i,t}$. Sau mỗi bước huấn luyện, nếu chuyên gia $i$ quá tải thì $b_i \leftarrow b_i - \gamma$, nếu thiếu tải thì $b_i \leftarrow b_i + \gamma$. Cân bằng tải được điều khiển trực tiếp như một bộ điều khiển phản hồi, không đi qua gradient.

\begin{center}
\bookimage[0.9\textwidth]{deepseek_v3_architecture}{Sơ đồ kiến trúc cơ bản của DeepSeek-V3: một khối Transformer gồm Multi-Head Latent Attention (MLA) với nén key-value hạng thấp và RoPE tách rời, cùng lớp DeepSeekMoE với chuyên gia chia sẻ (shared experts) và chuyên gia định tuyến (routed experts) được chọn bởi router top-K (Hình 2 trong DeepSeek-V3 Technical Report).}
\captionof{figure}{Kiến trúc khối Transformer của DeepSeek-V3: MLA cho attention và DeepSeekMoE cho FFN. Nguồn: DeepSeek-AI~\cite{deepseekai2024v3}.}
\label{fig:deepseek_v3_architecture}
\end{center}

\paragraph{Giải pháp bổ sung: Noisy Top-K Gating}
Một kỹ thuật khác là thêm nhiễu ngẫu nhiên (Gaussian noise) vào logits của mạng cổng trong quá trình huấn luyện.
$$ \text{logits}_{\text{noisy}} = \text{logits}_{\text{original}} + \text{Noise} $$
Nhiễu này buộc mô hình phải “khám phá” (explore) các lựa chọn định tuyến khác nhau thay vì chỉ đi theo con đường quen thuộc, giúp cải thiện sự cân bằng tải một cách tự nhiên.

\paragraph{2. Thách thức về Triển khai và Cơ sở hạ tầng}
\begin{itemize}
    \item \textbf{Bộ nhớ:} Mặc dù chi phí tính toán trên mỗi token thấp, toàn bộ các tham số của tất cả các chuyên gia vẫn phải được tải vào bộ nhớ của GPU, đòi hỏi các hệ thống có bộ nhớ cực lớn.
    \item \textbf{Giao tiếp mạng:} Trong các hệ thống huấn luyện phân tán, các token từ một GPU có thể cần được gửi đến các chuyên gia nằm trên một GPU khác, đòi hỏi băng thông mạng rất cao.
\end{itemize}

\subsection{Tổng kết về MoE}
Mixture-of-Experts là một sự thay đổi mô hình trong cách chúng ta xây dựng các mô hình ngôn ngữ lớn. Nó cho phép chúng ta tách rời \textbf{tổng số tham số} khỏi \textbf{chi phí tính toán trên mỗi token}.

Bằng cách sử dụng các chuyên gia thưa thớt, MoE mở ra một con đường để xây dựng các mô hình có quy mô khổng lồ (hàng nghìn tỷ tham số), có khả năng lưu trữ một lượng kiến thức cực lớn, trong khi vẫn duy trì được tốc độ huấn luyện và suy luận ở mức có thể quản lý được. Đến năm 2025--2026, MoE đã trở thành kiến trúc mặc định của hầu hết các mô hình mở hàng đầu (DeepSeek-V3, Qwen3-MoE, Llama 4, Kimi K2, gpt-oss), với xu hướng rõ rệt: \textbf{nhiều chuyên gia nhỏ hơn, tỷ lệ kích hoạt thấp hơn} (ví dụ 32B/1T ở Kimi K2), cùng với chuyên gia chia sẻ và các cơ chế cân bằng tải không dùng loss phụ.

\begin{tcolorbox}[title={Đọc paper MoE mới: bốn câu hỏi nên đặt ra}, colback=yellow!10!white, colframe=yellow!50!black, fonttitle=\bfseries]
\begin{enumerate}
	\item \textbf{Tổng tham số và tham số kích hoạt} là bao nhiêu? So sánh công bằng phải dựa trên FLOPs/token và bộ nhớ, không chỉ tổng tham số.
	\item \textbf{Cân bằng tải} bằng cách nào (loss phụ, bias, capacity) và có token nào bị bỏ (dropped) không?
	\item \textbf{Chi phí giao tiếp} được kiểm soát ra sao (giới hạn số node, expert parallelism)?
	\item Kết quả có \textbf{ablation cùng lượng compute} với mô hình dày đặc tương ứng không?
\end{enumerate}
\end{tcolorbox}